\section{Background}
\label{sec:background}

\subsection{CBR and Process-Oriented CBR}

Case-Based Reasoning \cite{aamodt1994,kolodner1993} solves new problems by retrieving
and adapting past cases. The four-stage cycle --- \textbf{Retrieve}, \textbf{Reuse},
\textbf{Revise}, \textbf{Retain} --- depends critically on case integrity: each case
must be a self-contained record of a past solution state, free of implicit context.

Process-Oriented CBR (POCBR) \cite{bergmann2002,minor2014} extends this to
\textbf{processes as cases}: the case captures a workflow trace including intermediate
states and data flow. Multi-step LLM workflows are a natural POCBR domain; the
challenge is ensuring their intermediate states are suitable, self-contained cases.

\subsection{Case Integrity and Structural Isolation}

Most LLM orchestration frameworks do not satisfy case integrity by default: a
checkpoint at step $k$ may embed accumulated conversation history or prompt context
from steps $1, \ldots, k-1$, making it a snapshot of shared state rather than a
self-contained case. Three consequences follow: \textbf{unreliable retrieval} ---
continuing from such a checkpoint in a new context produces incorrect behavior;
\textbf{non-localizable failures} --- the cause of a failure at step $k$ may be
anywhere in accumulated implicit state; \textbf{irreproducible cases} --- two
executions loading the same checkpoint may produce different outputs.

These failures are eliminated by \textbf{structural isolation}: a workflow language
provides structural isolation if each step can access only data explicitly declared in
its own scope block, the compiler determines each step's scope at formalization, and the
orchestrator constructs each step's input from only its declared references. Under this
guarantee, the stored output at step $k$ contains exactly what downstream steps need;
failures are localizable to the declared inputs of the failing step; and checkpoint
reuse reproduces the original data-dependency behavior. NormCode \cite{guan2026normcodesemiformallanguageauditable} is
designed around this enabling condition.
